\section{Introduction}\label{sec:1}

Online hiring markets route high application volume through platforms that rank candidates and jobs at scale, yet both sides often receive limited feedback on why a match appeared, changed, or disappeared.
Candidates submit the same résumé to many openings and cannot distinguish rejection, filtering, and queue effects.
When outcomes arrive---an interview invitation, a rejection, or a role removed from a results page---the rationale seldom returns to the applicant in actionable form.

Recruitment platforms centralize workflow, but they rarely make ranking decisions auditable.
Candidates see ordered lists without evidence tying each slot to profile fields or job requirements.
Recruiters and hiring managers see shortlists update after job-text edits without a visible link between requirement changes and rank movement.
Neither side can reliably answer: \emph{why was this person shown this role, and what would need to change for the ranking to differ?}

These gaps impose measurable cost.
Misaligned interviews consume calendar time on both sides; qualified candidates who never reach a shortlist extend vacancies or divert talent elsewhere.
Organizations pay in recruiter hours and manager attention; job seekers pay in repeated preparation for roles that were poor fits from the outset.

Assistive artificial intelligence can address such opacity without automating hire-or-reject decisions.
Sal Khan~\cite{Khan2024} argues that AI should guide people through complex choices---clarifying trade-offs, supporting preparation, and preserving human judgment.
In admissions and hiring alike, static snapshots often replace evolving context; software that maintains structured, versioned profiles can instead support informed decisions over time.

We instantiate that principle with a dedicated agent on each market side.
A \textbf{candidate-side agent}---also referred to as the client-side agent in our implementation---accepts résumés and profile updates, extracts a structured record of skills and experience, and keeps that record current as the user revises materials.
An \textbf{employer-side agent} performs the parallel work for hiring: it accepts job descriptions and requirements, organizes them into a structured job record, and updates that record when a role changes.
Each agent owns its side of the market; neither manages the other's data.

Separate ownership only helps if the two sides can still meet on common ground.
When candidate records and job records never connect, both portals collapse into blind browsing.
When one large program controls every edit and every score, users cannot see who changed what or why a recommendation moved.
JobMatch therefore treats \textbf{collaboration among agents} as the core design problem.
Candidate and employer agents maintain their own state; a neutral \textbf{Matchmaking Agent} reads both, compares them, and returns ranked recommendations with structured explanations.
When either side updates a résumé or a job description, the owning agent refreshes its record and notifies the Matchmaking layer so stale rankings can be invalidated.

These components form \textbf{JobMatch}, a multi-agent job-matching platform implemented as a research prototype with separate web portals for candidates, employers, and administrators.
The platform parses résumés and job descriptions into structured states, builds comparable representations of people and roles, and routes matching through the Matchmaking Agent to produce ordered recommendations.
For deployment scenarios, the Employer agent can optionally refresh its job corpus from an external provider API; offline evaluation in this paper uses a fixed demo corpus (\S\ref{sec:4}).
Users remain in control throughout: they review extracted fields, read why a pairing ranked as it did, and decide whether to save, apply, or shortlist.
The system advises; it does not hire or reject on anyone's behalf.
We evaluate the matching core offline on a labeled demo corpus---30 résumés, 15 jobs, and 47 graded relevance pairs---using standard ranking metrics and reproducible benchmark comparisons (\S\ref{sec:5}--\S\ref{sec:6}).
Hybrid rankers reach nDCG@5 up to $0.969$ on this corpus; the portal defaults to a fixed composite score whose components map to user-visible explanations.

This paper makes the following contributions:
\begin{itemize}[noitemsep]
  \item A \textbf{Candidate/Client Agent} that ingests résumés, maintains versioned candidate state, and updates candidate embeddings on confirmed saves.
  \item An \textbf{Employer Agent} that ingests job descriptions, maintains versioned job state, and optionally refreshes postings from an external provider API.
  \item A \textbf{Matchmaking Agent} that scores candidate--job pairs read-only and returns ranked lists with structured explanations.
  \item An \textbf{event-driven communication model} in which profile updates invalidate stale rankings without cross-agent writes to foreign stores.
  \item A \textbf{role-separated portal layer} demonstrating end-to-end candidate, employer, and administrator workflows.
  \item An \textbf{offline evaluation} on a fixed labeled corpus with benchmark comparisons, ablations, and a continuous-integration regression gate.
\end{itemize}

The remainder of the paper is organized as follows.
Section~\ref{sec:2} reviews recruitment systems, AI agents in hiring, and related matching approaches.
Section~\ref{sec:3} presents the multi-agent architecture (Candidate, Employer, and Matchmaking agents; communication; overall workflow).
Section~\ref{sec:4} describes implementation details (libraries, APIs, data structures, and parameters).
Section~\ref{sec:5} defines quality metrics, similarity scores, and evaluation criteria.
Section~\ref{sec:6} reports benchmark results and discussion.
Section~\ref{sec:7} concludes and outlines future scope.
